\begin{terminologyList}
    \term[Атака по сторонним каналам (Side-Channel Attack)]{класс криптографических атак, использующих информацию о физическом процессе выполнения алгоритма (например, время выполнения, энергопотребление), а не слабости самого алгоритма}
    \term[Векторизация (Vectorization)]{процесс преобразования алгоритма, оперирующего одиночными значениями (скалярами), в алгоритм, одновременно обрабатывающий массив данных с помощью SIMD-инструкций}
    \term[Конечные поля (Finite Fields)]{алгебраическая структура, состоящая из конечного множества элементов, для которых определены операции сложения, вычитания, умножения и деления}
    \term[Пакетная обработка (Batching)]{метод группировки нескольких независимых криптографических операций для их параллельного выполнения на одном процессорном ядре}
    \term[Проективные координаты (Projective Coordinates)]{система координат, позволяющая представлять точки эллиптической кривой в виде дробей, что исключает необходимость выполнения ресурсоемкой операции деления на каждом шаге сложения}
    \term[Сборщик мусора (GC, Garbage Collector)]{компонент платформы .NET, обеспечивающий автоматическое управление памятью и освобождающий память, занятую неиспользуемыми объектами}
    \term[AVX2 (Advanced Vector Extensions 2)]{набор SIMD-инструкций x86, позволяющий выполнять математические операции над 256-битными регистрами}
    \term[Bouncy Castle]{популярная кроссплатформенная криптографическая библиотека, де-факто являющаяся стандартом для многих управляемых (managed) сред}
    \term[Constant-time (Константное время выполнения)]{свойство программной реализации криптографического алгоритма, при котором время его работы не зависит от значений секретных ключей или обрабатываемых данных}
    \term[Curve25519 / X25519]{эллиптическая кривая в форме Монтгомери, используемая преимущественно для протоколов алгоритма Диффи-Хеллмана (ECDH)}
    \term[Ed25519 / EdDSA]{схема электронной цифровой подписи на основе искривленных эллиптических кривых Эдвардса, использующая хеш-функцию SHA-512}
    \term[Intrinsics (Интринсики)]{специальные методы в платформе .NET (\texttt{System.Runtime.Intrinsics}), транслируемые JIT-компилятором напрямую в специфичные машинные инструкции процессора}
    \term[Radix-51]{представление 255-битного числа в виде массива из пяти 64-битных элементов (лимбов), в каждом из которых значащими являются только младшие 51 бит}
    \term[SIMD (Single Instruction, Multiple Data)]{архитектурный принцип построения процессоров, позволяющий выполнять одну машинную инструкцию над набором данных одновременно}
    \term[Zero Allocation (Нулевые аллокации)]{подход к программированию в .NET, при котором в ходе выполнения критического пути кода не выделяется память в управляемой куче (heap), что исключает паузы на работу сборщика мусора}
    \term[.NET (.NET 10)]{кроссплатформенная среда выполнения и фреймворк разработки от Microsoft 10 версии, поддерживающая множество языков программирования и предоставляющая обширную стандартную библиотеку}

\end{terminologyList}